% !TEX root = ../Tesis.tex
\chapter{Metodología}
\section{Recolección de muestras}
El centro de muestreo se instaló en el Instituto Salamanca
($20.567158$, $-101.186115$) en el municipio de Salamanca, Guanajuato
(Figura~\ref{fig:mapa_sitio}). Las campañas de recolección se
desarrollaron del 20 de febrero al 5 de marzo de 2026 y del 25 al 29 de
mayo de 2026 mediante un sistema de condensación. El sistema está
compuesto por un refrigerador modificado marca Hisense, con una
conexión de tubería de PVC en su interior. El sistema tomaba las
muestras en un periodo entre las 22:00 y 2:00 h de forma automática.

% AGREGAR MAPA DE SITIO
\begin{figure}[H]
    \centering
    \includegraphics[width=0.85\textwidth]{figuras/mapa_sitio.png}
    \caption{Localización del sitio de muestreo en el Instituto
    Salamanca, Guanajuato.}
    \label{fig:mapa_sitio}
\end{figure}

\section{Análisis de $\delta^{2}$H y $\delta^{18}$O}
\label{sec:analisis_isotopico}

Para obtener un factor del fraccionamiento isotópico que se produce al
condensar la muestra, se adapta el condensador a una cámara de
ambientes controlados. Cada 24 horas se recolecta la muestra de vapor y
una muestra líquida. La cabina mantenía las condiciones en $30$\,°C y
$30\%$ de humedad relativa.

\section{Validación de datos de las estaciones de monitoreo de aire}
\label{sec:validacion_datos}

Los datos de las estaciones de monitoreo de aire recibieron un
tratamiento según el Protocolo de Manejo de Datos de Calidad del Aire
propuesto por el Instituto Nacional de Ecología (Instituto Nacional de
Ecología, s.f.). Se realiza una limpieza de los datos, asignando
banderas según los criterios de la Tabla~\ref{tab:banderas_limpieza}.

\begin{table}[H]
\centering
\footnotesize
\caption{Criterios de asignación de banderas para la limpieza de datos.}
\label{tab:banderas_limpieza}
\begin{tabular}{lp{11cm}}
\toprule
\textbf{Bandera} & \textbf{Criterio} \\
\midrule
VA & Dato válido: dato que se considera correcto o verdadero pero que
está sujeto a revisión en la etapa de verificación de datos. \\
IC & Dato inválido por calibración: datos que provengan de las
actividades de calibración de los equipos de medición. Asimismo, se
incluirán aquellos datos generados como resultado de la ejecución del
programa de pruebas de cero/span. \\
IR & Dato inválido por rango de operación: valores fuera de los límites
superior e inferior fijados en el equipo de medición
(Tablas~\ref{tab:rangos_operacion} y~\ref{tab:limites_inferiores}). \\
VZ & Dato válido igualado al límite de detección o a cero: datos con
valores negativos que, excepto en el caso de las mediciones de
temperatura, deberán ser igualados a cero o al límite de detección
correspondiente para poderlos incluir en la base de datos revisada. \\
IF & Dato inválido por falla del equipo: datos de los que exista
evidencia en bitácora o en listas de verificación cuyos parámetros de
operación (temperaturas, flujos, presiones, voltajes) se encuentran
fuera de los valores normales de operación. \\
ND & Dato no disponible: se aplica a cualquier dato ausente, ya sea por
fallas en la transmisión-recepción o por cortes de energía. \\
\bottomrule
\end{tabular}
\end{table}

\begin{table}[H]
\centering
\footnotesize
\caption{Rangos de operación por contaminante.}
\label{tab:rangos_operacion}
\begin{tabular}{ll}
\toprule
\textbf{Contaminante} & \textbf{Rango de operación} \\
\midrule
O\textsubscript{3} & 0--500 ppb \\
SO\textsubscript{2} & 0--500 ppb \\
NO\textsubscript{2} & 0--500 ppb \\
NO & 0--500 ppb \\
NO\textsubscript{x} & 0--500 ppb \\
CO & 0--50 ppm \\
Hidrocarburos totales & 0--10 / 0--20 ppmc \\
Metanos & 0--10 / 0--20 ppmc \\
No metanos & 0--10 / 0--20 ppmc \\
Partículas PM10 & 0--500 / 0--1000 \si{\micro\gram\per\cubic\meter} \\
Partículas PM2.5 & 0--500 / 0--1000 \si{\micro\gram\per\cubic\meter} \\
Humedad relativa & 0--100\% \\
Presión barométrica & 500--760 mmHg \\
Precipitación pluvial & 0--10 mm \\
Radiación solar & 0--2000 \si{\watt\per\square\meter} \\
Radiación UV-A & 0--300 mW/m\textsuperscript{2} \\
                  & 0--2000 mW/m\textsuperscript{2} \\
Radiación UV-B & 0--5 MED/hr \\
                & 0--600 mW/m\textsuperscript{2} \\
Velocidad del viento & 0--50 m/s \\
Dirección del viento & 0--360° \\
Temperatura ambiental & $-50$ a $500$\,°C \\
\bottomrule
\end{tabular}
\end{table}

\begin{table}[H]
\centering
\footnotesize
\caption{Límites inferiores por contaminante.}
\label{tab:limites_inferiores}
\begin{tabular}{ll}
\toprule
\textbf{Contaminante} & \textbf{Mínimo inferior} \\
\midrule
O\textsubscript{3} & $-3$ ppb \\
SO\textsubscript{2} & $-3$ ppb \\
NO\textsubscript{2} & $-3$ ppb \\
NO & $-3$ ppb \\
NO\textsubscript{x} & $-6$ ppb \\
CO & $-0.4$ ppm \\
\bottomrule
\end{tabular}
\end{table}

Para la verificación de datos, se asignan nuevas banderas de acuerdo con
la Tabla~\ref{tab:banderas_verificacion}, según los siguientes
criterios:

\begin{table}[H]
\centering
\footnotesize
\caption{Criterios de asignación de banderas para la verificación de
datos.}
\label{tab:banderas_verificacion}
\begin{tabular}{lp{2.5cm}p{8.5cm}}
\toprule
\textbf{Bandera} & \textbf{Significado} & \textbf{Criterio} \\
\midrule
DS & Dato sospechoso &
Cuando existen valores constantes por más de tres horas consecutivas
de CO, NO\textsubscript{x}, NO\textsubscript{2}, NO,
O\textsubscript{3}, SO\textsubscript{2}, PM10 o PM2.5. \\
IO & Inválido por operador &
Cuando la suma de NO y NO\textsubscript{2} dividida entre
NO\textsubscript{x} se encuentra fuera del intervalo $(0.85, 1.15)$. \\
   & &
Cuando la relación de PM2.5 entre PM10 es mayor a $1.15$. \\
   & &
Para datos de radiación (total, UVA, UVB, UVC) cuando esta es diferente
a cero durante la noche. \\
   & &
Para datos de velocidad de viento cuando esta no varía en más de
$0.1$\,m/s en tres horas consecutivas o en más de $0.5$\,m/s en doce
horas consecutivas. \\
   & &
Cuando los datos de dirección de viento no varían en más de $1^\circ$
por más de tres horas consecutivas o en más de $10^\circ$ durante
dieciocho horas consecutivas. \\
   & &
Cuando los datos de temperatura cambian en más de $5$\,°C respecto a la
hora previa o no varían en más de $0.5$\,°C en doce horas consecutivas. \\
   & &
Para datos de presión barométrica cuando esta tiene cambios de más de
$0.75$\,mmHg en tres horas consecutivas. \\
VC & Dato válido calculado &
Se usa en aquellos datos que son obtenidos por medio de interpolación o
por el uso de series de tiempo. También es aplicable a aquellos que son
corregidos de acuerdo con los resultados de las verificaciones de cero
y \textit{span}. \\
VE & Dato válido con evento extraordinario &
Se aplica para datos cuya frecuencia y duración fueron registrados en
la bitácora general de operación. También cuando existe evidencia
documental de eventos que afectaron la medición, tales como incendios,
tolvaneras o actividades de construcción. \\
DV & Dato para verificar &
Bandera temporal que incluye datos que deberán ser confirmados en alguna
de las categorías anteriores. También se puede asignar esta bandera a
datos que se presentan fuera de los límites establecidos con base en
tendencias históricas de cada uno de los contaminantes medidos o
parámetros meteorológicos. \\
\bottomrule
\end{tabular}
\end{table}

Para evaluar los datos estadísticamente se utilizó un análisis de datos
atípicos a través del cálculo de la anomalía estandarizada ($z$). Este
parámetro se calcula restando la media muestral de los datos brutos $x$
y dividiendo el resultado por la desviación estándar muestral
correspondiente (Ecuación~\ref{eq:z_score}).

\begin{equation}
z = \frac{x - \bar{x}}{s_x}
\label{eq:z_score}
\end{equation}

Se considera dato atípico cuando el valor de $z$ supera $\pm 2$ y
$\pm 3$ desviaciones estándar. En el primer caso, se categorizan como
sospechosos, mientras que en el segundo caso se etiquetan como dato
extremo.

Asimismo, se realizó una prueba de valores extremos mediante el rango
intercuartílico (Ecuación~\ref{eq:iqr}).

\begin{equation}
\text{IQR} = Q_3 - Q_1
\label{eq:iqr}
\end{equation}

Se consideran datos atípicos aquellos que permanecen fuera del rango
definido en la Ecuación~\ref{eq:rango_atipicos}. En tal caso, se
etiquetaron como valores extremos.

\begin{equation}
Q_1 - 1.5\,\text{IQR} > X < Q_3 + 1.5\,\text{IQR}
\label{eq:rango_atipicos}
\end{equation}